% ===== PRÉAMBULE CANONIQUE — à coller VERBATIM en tête de CHAQUE .tex =====
\documentclass[11pt,a4paper]{article}
\usepackage[utf8]{inputenc}\usepackage[T1]{fontenc}\usepackage{lmodern}
\usepackage[french]{babel}
\usepackage{amsmath,amssymb,mathtools}\usepackage{icomma}
\usepackage[version=4]{mhchem}          % \ce{...} : équations & formules
\usepackage{siunitx}\sisetup{locale=FR} % \qty{}{}, \si{}, \ang{}
\usepackage{chemfig}                    % \chemfig{...} : structures développées
\usepackage{xcolor}\usepackage{colortbl}
\usepackage{tikz}\usepackage{pgfplots}\pgfplotsset{compat=1.18}
\usepackage[most]{tcolorbox}\usepackage{enumitem}\usepackage{array,booktabs}
\usepackage[a4paper,margin=2.2cm]{geometry}
\usetikzlibrary{arrows.meta,calc,angles,quotes,positioning,patterns,backgrounds,shapes.geometric}
\definecolor{primary}{RGB}{222,49,99}\definecolor{primarydark}{RGB}{150,22,55}
\definecolor{defcol}{RGB}{0,114,178}\definecolor{methcol}{RGB}{52,92,148}
\definecolor{excol}{RGB}{0,135,90}\definecolor{attncol}{RGB}{200,40,55}
\tcbset{boxbase/.style={enhanced,breakable,boxrule=0.9pt,arc=2.5pt,
  left=3mm,right=3mm,top=2.3mm,bottom=2.3mm,fonttitle=\bfseries,coltitle=white}}
% --- boîtes TUTORIEL (noms EXACTS, ne pas renommer) ---
\newtcolorbox{cle}[1][]{boxbase,colback=defcol!6,colframe=defcol,
  title={\textsc{À retenir}\ifx\relax#1\relax\relax\else\textmd{ : \itshape #1}\fi}}
\newtcolorbox{methode}[1][]{boxbase,colback=methcol!7,colframe=methcol,sharp corners,
  title={\textsc{Méthode}\ifx\relax#1\relax\relax\else\textmd{ --- \itshape #1}\fi}}
\newtcolorbox{exemplebox}[1][]{boxbase,colback=excol!5,colframe=excol,
  title={\textsc{Exemple}\ifx\relax#1\relax\relax\else\textmd{ : \itshape #1}\fi}}
\newtcolorbox{piege}[1][]{boxbase,colback=attncol!6,colframe=attncol,
  title={\textsc{Piège}\ifx\relax#1\relax\relax\else\textmd{ : \itshape #1}\fi}}
\newtcolorbox{remarque}[1][]{boxbase,colback=black!4,colframe=black!45,
  title={\textsc{Remarque}\ifx\relax#1\relax\relax\else\textmd{ : \itshape #1}\fi}}
% --- boîtes EXERCICE / CORRIGÉ (noms EXACTS, auto-comptées) ---
\newtcolorbox[auto counter]{exo}[1][]{boxbase,colback=primary!4,colframe=primary,
  title={\textsc{Exercice}~\thetcbcounter\ifx\relax#1\relax\relax\else\textmd{ --- \itshape #1}\fi}}
\newtcolorbox[auto counter]{corrige}[1][]{boxbase,colback=excol!4,colframe=excol,
  title={\textsc{Corrigé}~\thetcbcounter\ifx\relax#1\relax\relax\else\textmd{ --- \itshape #1}\fi}}
\newcommand{\R}{\mathbb{R}}\newcommand{\N}{\mathbb{N}}\newcommand{\Z}{\mathbb{Z}}\newcommand{\ds}{\displaystyle}
% (un \usepackage{fancyhdr}+en-tête est possible ; il est ignoré côté web)
% =========================================================================
\begin{document}
\sisetup{per-mode=symbol}
\begin{corrige}[une pesée dans toutes les unités]
Changer d'unité ne fait que déplacer la virgule : les $3$ CS de la mantisse ($4,50$) sont conservés.
\begin{enumerate}[label=\alph*)]
  \item $\qty{1}{\gram}=10^{3}\ \si{\milli\gram}$ : $m = 4,50\times10^{0}\ \si{\milli\gram} = \qty{4,50}{\milli\gram}$.
  \item $\qty{1}{\gram}=10^{6}\ \si{\micro\gram}$ : $m = 4,50\times10^{3}\ \si{\micro\gram}$.
  \item $\qty{1}{\gram}=10^{-3}\ \si{\kilo\gram}$ : $m = 4,50\times10^{-6}\ \si{\kilo\gram}$.
\end{enumerate}
Les quatre écritures ont toutes $3$ CS.
\end{corrige}

\begin{corrige}[un même calcul, trois précisions]
On garde le nombre de CS voulu et l'on regarde le premier chiffre supprimé de $15,625$.
\begin{enumerate}[label=\alph*)]
  \item $4$ CS : chiffre supprimé $5\ge5 \Rightarrow \boxed{15,63}$ (par excès).
  \item $3$ CS : on garde $15,6$ ; chiffre supprimé $2<5 \Rightarrow \boxed{15,6}$.
  \item $2$ CS : on garde $15$ ; chiffre supprimé $6\ge5 \Rightarrow \boxed{16}$.
\end{enumerate}
\end{corrige}

\begin{corrige}[énergie de combustion par kilogramme]
Comme $\qty{1}{\gram}=10^{-3}\ \si{\kilo\gram}$ (facteur exact) :
\[ \frac{4,307\times10^{4}\ \si{\joule}}{\qty{1}{\gram}}
   =\frac{4,307\times10^{4}\ \si{\joule}}{10^{-3}\ \si{\kilo\gram}}
   =4,307\times10^{7}\ \si{\joule\per\kilo\gram}. \]
Arrondi à $3$ CS (chiffre supprimé $7\ge5$, par excès) :
$\boxed{\approx 4,31\times10^{7}\ \si{\joule\per\kilo\gram}}$.
\end{corrige}

\begin{corrige}[vrai ou faux, justifié]
\begin{enumerate}[label=\alph*)]
  \item \textbf{Vrai.} On garde $0,79$ ; le chiffre supprimé $9\ge5$ fait remonter une retenue :
        $0,79\to 0,80$ (les deux CS sont bien affichés).
  \item \textbf{Vrai.} $2,5\times10^{-3}=0,0025$ : c'est le même nombre, seule l'écriture change.
  \item \textbf{Faux.} La mantisse doit vérifier $1\le|a|<10$ : $12,5$ est interdit, on écrit $1,25\times10^{1}$.
  \item \textbf{Faux.} Arrondir trop tôt propage l'erreur : on ne doit arrondir qu'\emph{une seule fois},
        à la fin, en gardant des chiffres de garde aux étapes intermédiaires.
\end{enumerate}
\end{corrige}

\end{document}
